\documentclass[11pt,a4paper]{report}
\usepackage[utf8]{inputenc}
\usepackage[margin=1in]{geometry}
\usepackage{xcolor}
\usepackage{fancyhdr}
\usepackage{titlesec}
\usepackage{booktabs}
\usepackage{tabularx}
\usepackage{hyperref}
\usepackage{tcolorbox}
\usepackage{enumitem}
\usepackage{listings}
\usepackage{amsmath,amssymb}
\usepackage{graphicx}

% Colors
\definecolor{cdacblue}{RGB}{14, 53, 101}
\definecolor{labcyan}{RGB}{0, 150, 190}
\definecolor{darkgray}{RGB}{45, 52, 54}
\definecolor{lightgray}{RGB}{245, 247, 250}
\definecolor{alertred}{RGB}{214, 48, 49}
\definecolor{succgreen}{RGB}{46, 204, 113}

% Hyperref setup
\hypersetup{
    colorlinks=true,
    linkcolor=cdacblue,
    filecolor=cdacblue,      
    urlcolor=labcyan,
    citecolor=cdacblue,
    pdftitle={TrusMeet: PI and Administrator Operations Manual},
    pdfauthor={Krishna Saha, Dr. Dhiman Saha - de.ci.phe.red LAB, IIT Bhilai}
}

\lstdefinelanguage{json}{
    basicstyle=\footnotesize\ttfamily,
    breaklines=true,
    showstringspaces=false,
    frame=lines,
    backgroundcolor=\color{lightgray}
}
\lstset{
    backgroundcolor=\color{lightgray},
    basicstyle=\footnotesize\ttfamily,
    breaklines=true,
    captionpos=b,
    frame=lines,
    rulecolor=\color{cdacblue},
    keywordstyle=\color{cdacblue}\bfseries,
    stringstyle=\color{darkgray},
    commentstyle=\color{gray}\itshape,
    showstringspaces=false
}

% Page style
\pagestyle{fancy}
\fancyhf{}
\fancyhead[L]{\small\textcolor{cdacblue}{\textbf{TrusMeet}} --- PI \& Administrator Manual}
\fancyhead[R]{\small\textcolor{darkgray}{de.ci.phe.red LAB | IIT Bhilai}}
\fancyfoot[C]{\thepage}
\renewcommand{\headrulewidth}{0.5pt}

% Title formats
\titleformat{\chapter}[display]
  {\normalfont\huge\bfseries\color{cdacblue}}
  {\chaptertitlename\ \thechapter}{10pt}{\Huge}
\titleformat{\section}
  {\normalfont\Large\bfseries\color{cdacblue}}
  {\thesection}{1em}{}
\titleformat{\subsection}
  {\normalfont\large\bfseries\color{darkgray}}
  {\thesubsection}{1em}{}

\tcbuselibrary{skins,breakable}
\newtcolorbox{alertbox}[2][]{
  enhanced,
  breakable,
  colback=white,
  colframe=alertred,
  arc=3mm,
  boxrule=1pt,
  fonttitle=\bfseries\color{white},
  colbacktitle=alertred,
  title=#2,
  #1
}

\newtcolorbox{infobox}[2][]{
  enhanced,
  breakable,
  colback=white,
  colframe=cdacblue,
  arc=3mm,
  boxrule=1pt,
  fonttitle=\bfseries\color{white},
  colbacktitle=cdacblue,
  title=#2,
  #1
}

\newtcolorbox{tipbox}[2][]{
  enhanced,
  breakable,
  colback=white,
  colframe=succgreen,
  arc=3mm,
  boxrule=1pt,
  fonttitle=\bfseries\color{white},
  colbacktitle=succgreen,
  title=#2,
  #1
}

\begin{document}

% Title Page
\begin{titlepage}
    \centering
    \vspace*{1.5cm}
    
    {\Huge\bfseries\color{cdacblue} TrusMeet: Zero-Trust Lab Video Conference Platform\par}
    \vspace{0.8cm}
    {\LARGE\bfseries\color{darkgray} System Administrator \& Principal Investigator (PI) Manual\par}
    \vspace{0.4cm}
    {\large Complete Operations, Hardware PKCS\#11 Integration, and Cryptographic PDF Signing Guide\par}
    
    \vspace{2.5cm}
    \begin{tcolorbox}[colback=lightgray,colframe=cdacblue,width=0.9\textwidth,arc=2mm]
    \centering
    \textbf{Project:} Indigenous Cryptographic Hardware Token Integration in Academic Video Conferencing\\
    \textbf{Hardware Standard:} BIS IS-16695 (Parts 1 \& 2), ISO-7816, PKCS\#11 v2.20\\
    \textbf{Document Signing:} ISO 32000-1 / RFC 5652 CMS Detached Signatures (PAdES)
    \end{tcolorbox}
    
    \vspace{2.5cm}
    \begin{tabular}{ll}
    \textbf{Prepared By:} & Krishna Saha (AI Security) \\
    \textbf{Laboratory:} & de.ci.phe.red LAB (Cryptography \& Information Security) \\
    \textbf{Affiliation:} & Department of CSE, Indian Institute of Technology Bhilai \\
    \textbf{Principal Investigator:} & Dr. Dhiman Saha \\
    \textbf{Hardware Partner:} & C-DAC Noida \& IIT Bhilai \\
    \textbf{Version:} & 1.0 (Lab Edition) \quad | \quad \textbf{Date:} November 2026
    \end{tabular}
    
    \vfill
\end{titlepage}

\tableofcontents
\newpage

% CHAPTER 1
\chapter{System Overview \& Architecture}

\section{Introduction}
\textbf{TrusMeet} is an indigenous zero-trust video conferencing platform specifically engineered for research laboratory management, doctoral defense sessions, and confidential academic meetings. Unlike conventional commercial solutions (e.g., Zoom, Google Meet), TrusMeet enforces hardware-rooted identity validation via the **C-DAC TrusToken** cryptographic USB hardware device.

\section{Core Architectural Principles}
\begin{enumerate}[label=\textbf{\arabic*.}]
    \item \textbf{Hardware-Bound Identity}: Authentication is anchored in physical silicon. The user's private key never leaves the on-chip secure microcontroller.
    \item \textbf{Zero-Trust Admission}: Room gates can be configured to require strict hardware token validation. Non-token visitors cannot enter protected lab deliberations.
    \item \textbf{Continuous Presence Verification (Liveness)}: Authentication is not a one-time event; a periodic hardware heartbeat verifies token presence. Token removal triggers immediate security lockdowns.
    \item \textbf{Cryptographic Minutes Non-Repudiation}: Meeting resolutions and research agreements are digitally signed inside the PDF container using RFC 5652 CMS/PKCS\#7 structures.
\end{enumerate}

\section{System Component Stack}
\begin{table}[h!]
\centering
\small
\begin{tabularx}{\textwidth}{lX}
\toprule
\textbf{Subsystem} & \textbf{Role and Implementation Technologies} \\
\midrule
\textbf{WebRTC Signaling Server} & Node.js, Express, Socket.IO. Handles real-time mesh signaling, active room state, and audio/video transport coordination. \\
\textbf{PKCS\#11 Python Bridge} & \texttt{trustoken\_bridge.py} executing against \texttt{libtrustokenP11Libv3.3.0.so} and \texttt{libpcsclite.so.1}. Manages smartcard slots, PIN checks, and challenge signing. \\
\textbf{Cryptographic PDF Signer} & Standalone ISO 32000-1 / RFC 5652 engine (\texttt{pdf\_signer/}). Injects detached CMS signatures with cryptographic timestamps into meeting minutes. \\
\textbf{Edge Tunneling Service} & Cloudflare Edge routing via \texttt{start-public-tunnel.sh}. Provides trusted TLS termination for remote lab scholars without firewall port forwarding. \\
\bottomrule
\end{tabularx}
\caption{TrusMeet Architectural Stack}
\end{table}

\newpage
% CHAPTER 2
\chapter{Hardware Setup \& PKCS\#11 Driver Stack}

\section{Physical C-DAC TrusToken Dongle}
The C-DAC TrusToken is an indigenous cryptographic USB token developed jointly by C-DAC Noida and IIT Bhilai. It adheres to:
\begin{itemize}[noitemsep]
    \item \textbf{BIS IS-16695 (Part 1 \& 2)}: Indian Standard for Cryptographic Tokens.
    \item \textbf{ISO-7816}: Standard smart card interface specifications.
    \item \textbf{PKCS\#11 v2.20}: Cryptographic Token Interface standard (Cryptoki).
\end{itemize}

\subsection{USB Bus Identification}
When inserted into a Linux workstation, the kernel enumerates the device on the USB bus. You can confirm device presence using:
\begin{lstlisting}[language=bash]
lsusb
\end{lstlisting}
The expected output includes:
\begin{lstlisting}[language=bash]
Bus 001 Device 005: ID 288d:0040 CDAC INDIA TrusToken
\end{lstlisting}

\section{Understanding the Hardware LED Status}
The TrusToken hardware dongle is equipped with a status LED indicator:
\begin{itemize}
    \item \textbf{Solid Red Light (Normal / Standby)}: Confirms that the device is receiving 5V DC power from the host USB port, the microcontroller has initialized, and the crypto-container is in its default secure standby state awaiting PIN authentication.
    \item \textbf{Blinking / Flashing Red Light (Active Execution)}: Indicates active APDU command exchange between the operating system and the cryptographic chip (e.g., certificate reading, challenge signing, or PIN verification).
\end{itemize}

\section{Configuring the PC/SC Smartcard Daemon (\texttt{pcscd})}
The C-DAC PKCS\#11 library (\texttt{/usr/local/lib/libtrustokenP11Libv3.3.0.so}) relies on \texttt{libpcsclite} to communicate with CCID hardware via the PC/SC daemon.

\begin{alertbox}{Crucial Operational Prerequisite}
If \texttt{pcscd} is stopped, the PKCS\#11 bridge will fail with \texttt{CKR\_TOKEN\_NOT\_PRESENT (0xe0)} even if the USB token is physically inserted. You must start the daemon on the host.
\end{alertbox}

To start and enable the daemon on Ubuntu/Debian:
\begin{lstlisting}[language=bash]
sudo systemctl start pcscd
sudo systemctl enable pcscd
\end{lstlisting}

Verify daemon and reader status using the standard OpenSC utility:
\begin{lstlisting}[language=bash]
opensc-tool -l
\end{lstlisting}
You should see the active C-DAC TrusToken CCID reader listed.

\section{Validating Bridge Hardware Diagnostics}
To test token communication through the TrusMeet Python bridge, execute:
\begin{lstlisting}[language=bash]
cd /home/krishna-saha/trustoken-lab-meet
python3 trustoken_bridge.py hardware-status
\end{lstlisting}
When healthy, the response JSON will show:
\begin{lstlisting}[language=json]
{
  "library_present": true,
  "hardware_dongle_present": true,
  "slots_count": 1,
  "token_details": {
    "manufacturer": "CDAC-Noida (www.cdac.in)",
    "model": "TrusToken-v3",
    "serial_number": "CDAC-TT-2026-001"
  }
}
\end{lstlisting}

\newpage
% CHAPTER 3
\chapter{Server Deployment \& Network Orchestration}

\section{Starting the Application Server}
Navigate to the project root directory and start the Node.js application:
\begin{lstlisting}[language=bash]
cd /home/krishna-saha/trustoken-lab-meet
node server.js
\end{lstlisting}

The server binds to all interfaces (\texttt{0.0.0.0}) across two ports:
\begin{itemize}[noitemsep]
    \item \textbf{Port 8443 (HTTPS)}: Primary secure WebRTC port utilizing self-signed SSL certificates (\texttt{server.crt} and \texttt{server.key}).
    \item \textbf{Port 8080 (HTTP)}: Secondary port for local testing.
\end{itemize}

\section{WebRTC Browser Security: Why HTTPS is Mandatory}
Modern web browsers (Google Chrome, Mozilla Firefox, Microsoft Edge, Apple Safari) enforce the W3C \textbf{Secure Context} specification:
\begin{itemize}
    \item \texttt{http://localhost:8080} is specially grandfathered by browsers as secure.
    \item Any remote IP address accessed over plain HTTP (e.g., \texttt{http://10.10.1.226:8080}) is treated as \textbf{Insecure}. In this context, \texttt{navigator.mediaDevices.getUserMedia()} is completely disabled, blocking camera, microphone, and screen sharing.
\end{itemize}

\begin{tipbox}{Rule of Thumb for LAN Access}
When connecting over LAN IP, always use HTTPS: \texttt{https://<YOUR-IP>:8443}.
Accept the self-signed certificate warning by clicking \textbf{"Advanced" $\rightarrow$ "Proceed to IP (unsafe)"}.
\end{tipbox}

\section{Deploying the Cloudflare Public Mobile Tunnel}
To allow scholars and external collaborators on mobile phones or remote networks to join seamlessly without manual firewall routing or certificate errors, start the public tunnel:
\begin{lstlisting}[language=bash]
cd /home/krishna-saha/trustoken-lab-meet
./start-public-tunnel.sh
\end{lstlisting}
This script spawns \texttt{cloudflared} in the background and provisions a globally trusted HTTPS link (e.g., \texttt{https://dispatch-atmosphere-shed-title.trycloudflare.com/?room=LAB-DECIPHERED-SYNC}).

\section{Firewall Rules (UFW)}
If external clients on the local subnet cannot reach the server, ensure Ubuntu's firewall allows traffic:
\begin{lstlisting}[language=bash]
sudo ufw allow 8080/tcp
sudo ufw allow 8443/tcp
\end{lstlisting}

\newpage
% CHAPTER 4
\chapter{PI Host Privileges \& Zero-Trust Moderation}

\section{Member Onboarding \& Sign-Up Code Management}
To prevent unauthorized public internet users from creating accounts when using public edge tunnels, TrusMeet employs a \textbf{Secret Lab Registration Code}.

\subsection{Viewing and Managing the Sign-Up Code}
The code is stored on the server at:
\begin{lstlisting}[language=bash]
cat /home/krishna-saha/trustoken-lab-meet/data/signup-code.txt
\end{lstlisting}
The current active lab registration code is \texttt{\textbf{710b45f5787e}}.
\begin{itemize}
    \item \textbf{Custom Code}: The administrator can define a custom code by passing \texttt{SIGNUP\_CODE="MySecretCode"} before starting \texttt{node server.js}, or by editing \texttt{data/signup-code.txt}.
    \item \textbf{Timing Attack Resistance}: Code comparison is executed using constant-time SHA-256 digest comparison (\texttt{crypto.timingSafeEqual}) to prevent timing side-channel attacks.
    \item \textbf{Anti-Abuse Rate Limiting}: The server enforces strict rate limits: maximum 10 registration attempts per IP per hour, and 30 login attempts per 15 minutes.
\end{itemize}

\subsection{Cryptographic Account Issuance}
When a scholar registers with their name, username, and password:
\begin{enumerate}[noitemsep]
    \item An ECC P-256 / RSA keypair is minted via \texttt{trustoken\_bridge.py}.
    \item An X.509 certificate is signed under the Lab Root CA (\texttt{lab\_ca.crt}).
    \item The private key is encrypted with AES-256-GCM using a PBKDF2 derivative of the scholar's password, guaranteeing zero plaintext storage.
\end{enumerate}

\section{Host Role Differentiation}
The server evaluates user certificates issued under the Lab CA. A user logging in with the Principal Investigator certificate is automatically granted \textbf{Host Authority}.

\begin{table}[h!]
\centering
\small
\begin{tabular}{lcc}
\toprule
\textbf{Capability} & \textbf{PI (Host)} & \textbf{Lab Member (Scholar)} \\
\midrule
Audio/Video Stream Broadcast & \checkmark & \checkmark \\
Screen Share Presentations & \checkmark & \checkmark \\
In-Meeting Chat Participation & \checkmark & \checkmark \\
Room Locking / Unlocking & \checkmark & $\times$ \\
Forced Mute of Any Participant & \checkmark & $\times$ \\
Kick / Eject Participant & \checkmark & $\times$ \\
Official Sign-Off on Meeting Minutes & \checkmark & $\times$ (Review/Co-sign only) \\
\bottomrule
\end{tabular}
\caption{Role Permissions Matrix}
\end{table}

\section{Room Security Controls}
\begin{itemize}
    \item \textbf{Locking a Room}: The PI can toggle the room lock button in the top navigation bar. When locked, incoming connection requests from new scholars are rejected until unlocked.
    \item \textbf{Participant Moderation}: Clicking on a participant's tile provides options to mute noisy microphones or disconnect unauthorized users.
\end{itemize}

\section{The Physical Token Presence Heartbeat}
\begin{infobox}{How the Heartbeat Operates}
Every 12 seconds, the client sends a nonce challenge to the TrusToken bridge. If the USB dongle is removed from the USB port:
\begin{enumerate}[noitemsep]
    \item The PKCS\#11 bridge reports failure to sign the micro-challenge.
    \item The client immediately silences the user's microphone and pauses video transmission.
    \item A broadcast message (\texttt{peer-token-disconnected}) alerts all meeting attendees that the physical token has been removed.
\end{enumerate}
\end{infobox}

\newpage
% CHAPTER 5
\chapter{Cryptographic PDF Signing of Meeting Minutes}

\section{Standard-Compliant Cryptographic Protocol}
TrusMeet includes a complete PDF cryptographic signing engine located in \texttt{pdf\_signer/}. It strictly conforms to:
\begin{itemize}[noitemsep]
    \item \textbf{ISO 32000-1 / ISO 32000-2}: Standard PDF document architecture.
    \item \textbf{RFC 5652 (CMS / PKCS\#7)}: Cryptographic Message Syntax for detached signatures.
    \item \textbf{ETSI EN 319 142 (PAdES)}: PDF Advanced Electronic Signatures.
    \item \textbf{RFC 5280}: X.509 Certificate and CRL specifications.
\end{itemize}

\section{Signing Workflow via CLI}
To sign any lab document or generated meeting minutes with the PI's digital signature:
\begin{lstlisting}[language=bash]
cd /home/krishna-saha/trustoken-lab-meet/pdf_signer
python3 demo.py sign --input minutes.pdf --output minutes_signed.pdf --name "Dr. Dhiman Saha"
\end{lstlisting}

\subsection{The Automated Test and Verification Suite}
To execute the comprehensive end-to-end cryptographic test (PKI generation, visual stamp injection, detached CMS calculation, deep verification, and intentional tamper detection):
\begin{lstlisting}[language=bash]
python3 demo.py full-flow
\end{lstlisting}

\section{Universal Verification (``Verify Anywhere'')}
Signed PDFs produced by TrusMeet are verified by any standard compliant utility:
\begin{enumerate}
    \item \textbf{Linux CLI (Poppler \texttt{pdfsig})}:
\begin{lstlisting}[language=bash]
pdfsig minutes_signed.pdf
\end{lstlisting}
Output confirms: \textit{"Signature is Valid"} and identifies the signer certificate.
    \item \textbf{Adobe Acrobat Reader}:
    Open the signed PDF in Adobe Reader. The signature banner will report a valid digital signature adhering to PAdES-B.
\end{enumerate}

\section{Mathematical Tamper Resistance}
The signature covers the exact document byte range excluding the signature container itself (\texttt{/ByteRange} allocation). Modifying a single character or byte in the document invalidates the SHA-256 digest, immediately alerting any reader to unauthorized alteration.

\newpage
% CHAPTER 6
\chapter{Troubleshooting \& Diagnostic Reference}

\begin{table}[h!]
\centering
\small
\begin{tabularx}{\textwidth}{p{4cm}Xp{4.5cm}}
\toprule
\textbf{Symptom} & \textbf{Probable Root Cause} & \textbf{Resolution Step} \\
\midrule
\texttt{CKR\_TOKEN\_NOT\_PRESENT} error in logs. & USB dongle unplugged OR \texttt{pcscd} service is stopped. & Run \texttt{lsusb} to verify hardware, then execute \texttt{sudo systemctl restart pcscd}. \\
\midrule
Remote mobile browser cannot turn on camera/mic. & Page accessed via plain HTTP over network IP. & Connect using the HTTPS URL or start \texttt{./start-public-tunnel.sh}. \\
\midrule
Public tunnel fails to obtain URL. & \texttt{cloudflared} process already running or dead lockfile. & Run \texttt{pkill cloudflared} and restart \texttt{./start-public-tunnel.sh}. \\
\midrule
Browser displays ``Your connection is not private''. & Self-signed certificate on port 8443. & Click ``Advanced'' $\rightarrow$ ``Proceed to localhost (unsafe)''. \\
\midrule
PIN authentication rejected. & Incorrect hardware PIN entered. & TrusToken allows limited PIN retries before lockup. Re-verify the PIN before re-submitting. \\
\bottomrule
\end{tabularx}
\caption{Administrator Troubleshooting Matrix}
\end{table}

\vspace{2cm}
\begin{center}
\textbf{--- End of PI \& Administrator Operations Manual ---}\\
\small \textit{de.ci.phe.red LAB $\cdot$ Department of CSE $\cdot$ Indian Institute of Technology Bhilai}
\end{center}

\end{document}
